\subsection{The total restriction splitting}
\label{subsec:syntactic-total-splitting}

The syntax and theory construction of Appendix~\ref{app:syntax-theories}
will be applied in
\[
 \mathcal B_T=\Tot(\Split(S_T)).
\]
We describe this category by arithmetic predicates and prove the
required properties. Syntactic partial maps presented by functional
$\Sigma_1$ formulas and provable graph equality are discussed by
Robinson and Rosolini~\cite[pp.~108--109]{RobinsonRosolini}.

Let $\mathcal P_T$ have as objects the $\Sigma_1$ arithmetic formulas
$C(x)$ with the displayed free variable. A morphism from $C$ to $D$
is represented by a $\Sigma_1$ formula $\Gamma(x,y)$ for which $T$
proves
\begin{align}
 &\forall x\forall y\,
   \bigl(\Gamma(x,y)\longrightarrow C(x)\land D(y)\bigr),
                                               \label{eq:syntactic-support}\\
 &\forall x\,
   \bigl(C(x)\longrightarrow\exists!y\,\Gamma(x,y)\bigr).
                                               \label{eq:syntactic-totality}
\end{align}
Two representing formulas define the same morphism when $T$ proves
their universal equivalence. The identity on $C$ is represented by
$C(x)\land y=x$, and composition is represented by
\[
 (\Delta\Gamma)(x,z)
   \equiv\exists y\,(\Gamma(x,y)\land\Delta(y,z)).
\]
These formulas are $\Sigma_1$ and satisfy
\eqref{eq:syntactic-support}--\eqref{eq:syntactic-totality}; first-order
logic proves the identity and associativity equations. We write
$\Gamma:C\to D$ for the morphism represented by $\Gamma(x,y)$.
Provably equivalent carrier predicates need not be the same object
of this presentation, but the identity relations give an isomorphism
between them.

\begin{lemma}\label{lem:syntactic-functional-search}
For every $\Sigma_1$ formula $\Gamma(x,y)$ there is an index
$a_\Gamma$ such that PA proves
\begin{gather*}
 \forall x\forall y\,
   \bigl(G_{a_\Gamma}(x,y)\longrightarrow\Gamma(x,y)\bigr),\\
 \forall x\,
   \bigl(\exists y\,G_{a_\Gamma}(x,y)
          \longleftrightarrow\exists y\,\Gamma(x,y)\bigr).
\end{gather*}
If $T$ proves that $\Gamma$ is functional, then
\[
 T\vdash\forall x\forall y\,
       \bigl(G_{a_\Gamma}(x,y)\longleftrightarrow\Gamma(x,y)\bigr).
\]
\end{lemma}
\begin{proof}
Write $\Gamma(x,y)$ in the PA-provably equivalent form
$\exists w\,R(x,y,w)$, with $R$ primitive recursive. Search the
codes $n=0,1,\ldots$ in increasing order for
$R(x,\pi_0n,\pi_1n)$, and return $\pi_0n$ at the first successful
test. Here $\pi_0,\pi_1$ are the projections of the fixed numerical
pairing. The program and its index depend on $\Gamma$, not on $T$.

PA proves by induction on the number of tests that a halting run
returns the first successful test's first component. Conversely,
if there is a successful test, the least such code supplies a finite
halting computation. Formalized least-number search and the
computation coding therefore prove both displayed assertions.
If $T$ also proves functionality, any pair satisfying $\Gamma$
has the same output as this successful search. This gives the
last universal graph equivalence.
\end{proof}

\begin{lemma}\label{lem:syntactic-total-category}
There is an equivalence
\[
 \mathcal P_T\simeq\mathcal B_T.
\]
The category $\mathcal B_T$ is regular, with finite coproducts
that are disjoint and stable under pullback. Every monomorphism
of $\mathcal B_T$ has a partial inverse in $\Split(S_T)$.
For $\Gamma:C\to D$ in the predicate presentation, its regular
image is represented by the inclusion
\[
 R_\Gamma(y)\equiv\exists x\,\Gamma(x,y)
 \quad\hookrightarrow\quad D(y).
\]
In particular, every subobject of a predicate $D$ is isomorphic
over $D$ to an inclusion of a $\Sigma_1$ predicate.
\end{lemma}
\begin{proof}
Send a carrier $C$ to $(\omega,[C]_T)$ and a represented morphism
$\Gamma:C\to D$ to $[a_\Gamma]_T$. By
Lemma~\ref{lem:syntactic-functional-search}, its graph is
$T$-provably equivalent to $\Gamma$. Conditions
\eqref{eq:syntactic-support}--\eqref{eq:syntactic-totality} imply
\[
 [a_\Gamma]_T=[D]_T[a_\Gamma]_T[C]_T,
 \qquad\overline{[a_\Gamma]_T}=[C]_T.
\]
Thus it is a total morphism between the indicated split objects.
The arithmetized composition equation proves functoriality.

Conversely, a total morphism from $(\omega,[C]_T)$ to
$(\omega,[D]_T)$ has a graph satisfying
\eqref{eq:syntactic-support}--\eqref{eq:syntactic-totality} in $T$.
Equality of these morphisms is exactly universal $T$-provable graph
equivalence. The functor is therefore full and faithful.
For essential surjectivity, use the coding isomorphism $X\cong\omega$
to transport any object $(X,e)$ of $\Split(S_T)$ to $(\omega,e')$.
If $b$ represents $e'$, put $C(x)\equiv\exists y\,G_b(x,y)$.
The equation $e'=\bar e'$ says that $T$ proves
\[
 G_b(x,y)\longleftrightarrow(C(x)\land y=x)
\]
universally. Hence $e'=[C]_T$, proving the equivalence.

We prove the remaining claims in $\mathcal P_T$.
The predicate $x=0$ is terminal: the unique morphism to it from
$C$ has graph $C(x)\land y=0$.
A product of $C$ and $D$ is represented by
\[
 C(\pi_0z)\land D(\pi_1z),
\]
with its numerical projections. If $\Gamma:A\to C$ and
$\Delta:A\to D$, their pairing has graph
\[
 \exists x\exists y\,
   \bigl(\Gamma(a,x)\land\Delta(a,y)\land z=\langle x,y\rangle\bigr).
\]
The PA-provable pairing identities give its universal property.
For parallel morphisms $\Gamma,\Delta:C\to D$, the equalizer is
the inclusion into $C$ of
\[
 E(x)\equiv C(x)\land
             \exists y\,(\Gamma(x,y)\land\Delta(x,y)).
\]
Indeed, if $\Lambda:A\to C$ equalizes $\Gamma$ and $\Delta$,
equality of the composite graphs and functionality of $\Lambda$
give a $T$-proof that $\Lambda(a,x)$ implies $E(x)$.
The same graph therefore defines its unique factorization through
$E$. This proves finite completeness. In particular, the pullback
of $\Gamma:C\to E$ and $\Delta:D\to E$ is represented by
\[
 P(w)\equiv C(\pi_0w)\land D(\pi_1w)\land
       \exists y\,(\Gamma(\pi_0w,y)\land\Delta(\pi_1w,y)),
\]
with the two numerical projections.

Call a represented morphism $\varepsilon:C\to D$
\emph{provably surjective} when
\[
 T\vdash\forall y\,
       \bigl(D(y)\longrightarrow\exists x\,\varepsilon(x,y)\bigr).
\]
Every such morphism has a section. Apply the first two assertions
of Lemma~\ref{lem:syntactic-functional-search} to the relation
$\varepsilon(x,y)$ with input $y$ and output $x$. The resulting
program chooses the first coded witness of a preimage. Its graph
is PA-provably functional, even though the transposed relation
need not be functional. Provable surjectivity makes its domain
$D$, and the support of $\varepsilon$ makes its outputs belong
to $C$. It therefore defines $\sigma:D\to C$. The search
correctness statement and functionality of $\varepsilon$ give
$\varepsilon\sigma=1_D$.

Such a split epimorphism is the coequalizer of its kernel pair.
Indeed, let $k_0,k_1:K\rightrightarrows C$ be the kernel-pair
projections, and let $\beta:C\to Z$ satisfy
$\beta k_0=\beta k_1$. Since
$\varepsilon=\varepsilon\sigma\varepsilon$, the pair
$(1_C,\sigma\varepsilon)$ induces a morphism $C\to K$.
Thus $\beta=\beta\sigma\varepsilon$, so
$\beta\sigma:D\to Z$ is the required factorization. If
$\lambda\varepsilon=\beta$, then
$\lambda=\lambda\varepsilon\sigma=\beta\sigma$, proving
uniqueness.

Every $\Gamma:C\to D$ factors through $R_\Gamma$ by the same
graph, followed by the identity inclusion of $R_\Gamma$ into $D$.
The first morphism is provably surjective, hence regular epic by
the preceding argument, and the inclusion is monic.
Conversely, a regular epimorphism factors through its predicate
image as above; extremality of regular epimorphisms forces the
image inclusion to be an isomorphism, which implies its provable
surjectivity. Thus regular epimorphisms are precisely the provable
surjections, and they all split. A pullback of a split epimorphism
$\varepsilon:C\to D$ along $\Gamma:Z\to D$ has the section
$\langle\sigma\Gamma,1_Z\rangle:Z\to C\times_D Z$.
Consequently regular epimorphisms are stable under pullback.
This proves regularity and the asserted image formula.

The false predicate is initial. If $C$ admits a morphism to it,
totality gives $T\vdash\forall x\,\neg C(x)$, so $C$ is
isomorphic to the false predicate; thus the initial object is
strict. A coproduct is the tagged predicate
\[
 (C+D)(z)\equiv
   \exists x\,(z=2x\land C(x))\;\lor\;
   \exists y\,(z=2y+1\land D(y)),
\]
with the even and odd injections. Copairing is defined by selecting
the graph corresponding to the tag. The parity identities prove
the coproduct equations and disjointness of its injections.
For $\Gamma:A\to C+D$, put
\[
 A_C(a)\equiv\exists x\,\Gamma(a,2x),\qquad
 A_D(a)\equiv\exists y\,\Gamma(a,2y+1).
\]
These are the pullback carriers of the two injections, up to the
isomorphism that forgets the uniquely determined output coordinate.
Totality and functionality of $\Gamma$ prove in $T$ that
\[
 A(a)\longleftrightarrow A_C(a)\lor A_D(a),\qquad
 \neg(A_C(a)\land A_D(a)).
\]
The canonical morphism $A_C+A_D\to A$ forgets the tag.
Its inverse has graph
\[
 (A_C(a)\land z=2a)\;\lor\;(A_D(a)\land z=2a+1).
\]
The partition equations prove both inverse identities. Hence finite
coproducts are disjoint and stable under pullback.

It remains to analyze arbitrary monomorphisms. Let
$\Gamma:C\to D$ be monic and form its kernel pair. Its carrier is
\[
 K(w)\equiv C(\pi_0w)\land C(\pi_1w)\land
    \exists y\,(\Gamma(\pi_0w,y)\land\Gamma(\pi_1w,y)).
\]
Monicity implies equality of the two kernel-pair projections.
Consequently, $T$ proves $K(w)\to\pi_0w=\pi_1w$ universally.
Using the numerical pairing, this gives
\[
 T\vdash\forall x\forall x'\forall y\,
  \bigl(\Gamma(x,y)\land\Gamma(x',y)\longrightarrow x=x'\bigr).
\]
Regard the monomorphism in $\Split(S_T)$ through the preceding
equivalence. The transposed relation $\Gamma(x,y)$, with input $y$ and
output $x$, is therefore $T$-provably functional. By
Lemma~\ref{lem:syntactic-functional-search} it defines a partial
morphism $r:D\rightharpoonup C$ in $\Split(S_T)$ whose domain
predicate is $R_\Gamma$. Totality and functionality of $\Gamma$,
together with the displayed injectivity statement, prove
\[
 r\Gamma=1_C,\qquad
 \Gamma r=[R_\Gamma]_T=\bar r.
\]
Here $[R_\Gamma]_T$ denotes the partial identity on the predicate
$R_\Gamma$, regarded as a partial endomorphism of the split object
presented by $D$. Thus every monomorphism has a partial inverse.
Moreover, the corestriction $C\to R_\Gamma$ is an isomorphism,
with the same transposed relation as its total inverse. This proves
the assertion about subobjects and completes the proof.
\end{proof}

In particular, there is a restriction equivalence
\[
 \Par(\mathcal B_T,\Mono)\simeq\Split(S_T).
\]
Explicitly, a span $C\xleftarrow{m}E\xrightarrow{g}D$ is sent
to $gm^{-1}$, where $m^{-1}$ is the partial inverse just constructed.
Conversely, a partial morphism is represented by its domain
inclusion and its total restriction to that domain. In the predicate
presentation the domain is $\exists y\,G_f(x,y)$; relational
composition identifies composition of these spans by pullback with
composition in $\Split(S_T)$.

The category $\Split(S_T)$ is a Turing category,
with Turing object $U=(\omega,1_\omega)$.
Indeed, the displayed restriction equivalence supplies a restriction
terminal object, represented by the singleton predicate $x=0$.
Binary restriction products are
$(X,e)\times(Y,d)=(X\times Y,e\times d)$.
For every object $(X,e)$, the total morphism
$e:(X,e)\to(X,1_X)$ has retraction
$e:(X,1_X)\rightharpoonup(X,e)$, and $(X,1_X)\cong U$
by the coding isomorphism.
The Turing morphism extends by the idempotent-splitting argument
of Cockett--Hofstra~\cite[Section 3.5]{CockettHofstra}.
Explicitly, let $\tau:\omega\times\omega\rightharpoonup\omega$
be the Turing morphism of $S_T$, and let
$f:(X,e)\times U\rightharpoonup U$ be a morphism of $\Split(S_T)$.
Its representative $f:X\times\omega\rightharpoonup\omega$ in $S_T$
satisfies $f=f(e\times1_\omega)$. Choose a total morphism
$r:X\to\omega$ in $S_T$ with $f=\tau(r\times1_\omega)$.
Then $re:(X,e)\to U$ is total, since $\overline{re}=e$, and
\[
 f=\tau(r\times1_\omega)(e\times1_\omega)
   =\tau((re)\times1_\omega).
\]
The recognition criterion~\cite[Theorem 3.4]{CockettHofstra}
therefore applies.

Ranges extend to $\Split(S_T)$ by the same representatives.
For a morphism $f:(X,e)\rightharpoonup(Y,d)$, the equation
$df=f$ implies $\ran_{S_T}(f)\leq d$. This range is therefore
a restriction idempotent on $(Y,d)$, with the same leastness
and agreement properties. Lemma~\ref{lem:range-regular-image}
applies because every monomorphism of $\mathcal B_T$ has a
partial inverse.

\subsection{Syntax and lists}
\label{sec:st-parameterized-syntax}

We use primitive-recursive codes for finite sequences with length,
entry, and concatenation operations whose defining equations are
provable in $\PA$. Each finite sequence has one code; the empty
sequence and adjoining an entry are injective constructors with
disjoint images. For example, encode the empty sequence by $0$
and $(a,w)$ by $1+\langle a,w\rangle$. The tail code is smaller
than the nonempty sequence code, so bounded traversal defines the
sequence operations and induction in $\PA$ proves their equations.
All displayed arithmetical implications and equivalences in this
subsection are understood to be universally quantified.

\begin{lemma}[Finite membership conditions]
\label{lem:st-finite-sigma1}
Let $C(x)$ be $\Sigma_1$. The assertion that all entries of a coded
finite sequence satisfy $C$ is equivalent in $\mathrm{PA}$ to a
$\Sigma_1$ formula. The same holds for finitely many different
$\Sigma_1$ conditions selected by the primitive-recursive tags of
the entries.
\end{lemma}
\begin{proof}
Write $C(x)$ as $\exists v\,C_0(x,v)$ with $C_0$ primitive recursive.
For a sequence code $a$, use a coded sequence of witnesses:
\[
 \exists b\,\bigl(\operatorname{len}(b)=\operatorname{len}(a)
   \land\forall i<\operatorname{len}(a)\,
                  C_0(a_i,b_i)\bigr).
\]
Reading the entries proves the reverse implication. For the forward implication,
induction on the length constructs a sequence of witnesses by
appending one witness at a time. This is an instance of induction in $\PA$.
For finitely many tags, record the appropriate witness in each
position and use the corresponding finite case distinction.
\end{proof}

Fix a nonnegative integer $d$ and objects
\[
 V,\quad \Sigma_0,\ldots,\Sigma_d,\quad
 \Pi_0,\ldots,\Pi_d
\]
of $\mathcal B_T$, represented by $\Sigma_1$ carrier formulas with
the same letters. Numerical equality gives decidable equality on
$V$: on the source $V\times V$, the Boolean test is true precisely
when its two numerical coordinates are equal. The diagonal and its
complement are therefore represented by the predicates
$V(x)\land V(y)\land x=y$ and
$V(x)\land V(y)\land x\ne y$, respectively.

We choose codes for raw, finite, ordered syntax trees as follows.
A code is a canonical sequence of node records in postfix order.
A record contains one of the finitely many constructor tags and,
when the constructor has a label, that numeric label. Its ordered
child blocks are recovered by parsing the preceding records with
the arities specified by their tags. A term variable has arity zero; a
$k$-ary function or relation record has arity $k$; equality and
implication have arity two; negation and universal quantification
have arity one. A universal-quantifier record also contains its
variable label. A primitive-recursive stack parser checks the
sorts of the children and that exactly one tree remains.

Each constructor concatenates its ordered child blocks and appends
its root record. Its decomposition is unique. The number of records
is a primitive-recursive size function, and each child has smaller
size. All these statements are proved in $\mathrm{PA}$ by induction
on the length of the parsed sequence.

The formula $\operatorname{Tm}(t)$ asserts that this parser accepts
$t$ as a term, that variable labels satisfy $V$, and that each
function label of arity $k$ satisfies $\Sigma_k$. The formula
$\operatorname{Fm}(a)$ has the analogous conditions for formulas,
including membership in $\Pi_k$ for relation labels and in $V$ for
quantifier labels. Lemma~\ref{lem:st-finite-sigma1} makes these
formulas $\Sigma_1$: membership witnesses are existentially
quantified and are not components of a syntax code.
The constructor decomposition equivalences for these carrier
formulas are provable in $\mathrm{PA}$. For example, the universal
constructor has the equivalence
\[
 \operatorname{Fm}(\operatorname{all}(x,a))
       \longleftrightarrow V(x)\land\operatorname{Fm}(a),
\]
where $\operatorname{all}$ is the numerical universal-quantifier
constructor. The corresponding equivalences for a function or relation
constructor include its coefficient-carrier predicate and the
carrier predicates of its ordered children.
Write
\[
 \Tm=(\omega,[\operatorname{Tm}]_T),\qquad
 \Fm=(\omega,[\operatorname{Fm}]_T)
\]
for the corresponding objects of $\mathcal B_T$.

\begin{proposition}[Syntax initiality with parameters]
\label{prop:st-syntax-initiality}
Let $d$ be a nonnegative integer and let $V,\Sigma_k,\Pi_k$ be
objects of $\mathcal B_T$ for $0\leq k\leq d$. Powers below
denote finite products. The endofunctor
$K:\mathcal B_T^2\to\mathcal B_T^2$ defined by
\[
 K(X,Y)=\left(
 V+\coprod_{k=0}^{d}\Sigma_k\times X^k,
 X^2+\coprod_{k=0}^{d}\Pi_k\times X^k+Y+Y^2+V\times Y
 \right)
\]
has an initial algebra in $\mathcal B_T^2$, with carriers
$\Tm$ and $\Fm$ as constructed above. This algebra is stable
under parameters in the sense of Appendix~\ref{app:syntax-theories}: for every object $P$,
its pullback to $(\mathcal B_T/P)^2$ is initial for $K_P$.
\end{proposition}
\begin{proof}
The primitive-recursive constructors, restricted to their displayed
sources, are morphisms in $\mathcal B_T$ by the constructor
decomposition equivalences. They give the claimed algebra.

Fix an arbitrary object $P$ and an arbitrary $K_P$-algebra with
carriers $\xi:X\to P$ and $\upsilon:Y\to P$. Choose graph indices for their projections and for the
finitely many algebra operations. Write
$\Xi(x,p)$ and $\Upsilon(y,p)$ for the projection graphs. The source of a function operation is represented, for example, by
the predicate
\[
       P(p)\land\Sigma_k(s)\land
       \bigwedge_{i=1}^k\bigl(X(x_i)\land\Xi(x_i,p)\bigr).
\]
The algebra operation
is $T$-provably total on this predicate and has values $x$ satisfying
$X(x)\land\Xi(x,p)$. The other operations have the corresponding
fiber-product sources and preserve their projection to $P$.

Define relations $\operatorname{Ev}_{\mathrm{Tm}}(p,t,x)$ and
$\operatorname{Ev}_{\mathrm{Fm}}(p,a,y)$ by finite evaluation
records. A record assigns to every node its output value
in $X$ or $Y$, as appropriate, together with a halting computation
trace for the relevant algebra operation. At a node, the inputs
are its label, the values already assigned to its children, and
the fixed parameter $p$. Require the appropriate carrier and
projection conditions for every assigned value. The last assigned
value is the displayed output. Include $P(p)$ and the appropriate
syntax-carrier predicate in the relation.

All these requirements are $\Sigma_1$. Boundedly many carrier and
projection requirements are combined using
Lemma~\ref{lem:st-finite-sigma1}; the operation-trace requirements
are decidable checks on the recorded traces. Existentially
quantifying over the entire record therefore gives a
$\Sigma_1$ formula.

The following is an induction proof in $T$. Induct simultaneously
for terms and formulas on the number of nodes. Its assertion is
that every valid input of size at most $n$, with every parameter
satisfying $P$, has exactly one evaluation value, of the required
sort and over that parameter. At a variable or nullary constructor,
the corresponding algebra operation is total on the input
predicate. At any other constructor, its unique decomposition gives
children of smaller size. The induction hypothesis supplies their
values with the required common projection. They therefore form
an input in the source of the algebra operation at the root.
Its totality supplies an output in the correct fiber. Concatenating
the child records and appending the root computation supplies
the evaluation record. Functionality of the operation graphs,
together with the induction hypothesis, proves uniqueness of the
value. Evaluation records need not be unique.

This argument uses the finitely many totality and projection
equations of the given algebra as theorems of $T$. Its induction formula is an arithmetic formula, so the required
induction is available from $\PA$.

Lemma~\ref{lem:syntactic-functional-search} now supplies indices for
the two evaluation relations. The preceding proof verifies in $T$
that they are total morphisms
\[
 h_{\mathrm{Tm}}:P\times\Tm\longrightarrow X,
 \qquad
 h_{\mathrm{Fm}}:P\times\Fm\longrightarrow Y,
\]
and that they commute with the projections to $P$. Decomposing a
record at its root proves, in $T$, their algebra-morphism
equations: evaluation of a constructor is the corresponding algebra
operation applied to the evaluated children. Thus these arrows form
a $K_P$-algebra morphism in the slice.

Let $(h'_{\mathrm{Tm}},h'_{\mathrm{Fm}})$ be any other such algebra
morphism. Its totality and its algebra-morphism equations are
theorems of $T$. A second simultaneous induction on size proves
that its values equal the evaluation values. At a constructor,
use its algebra-morphism equation, replace the child values by the
induction hypothesis, and use functionality of the algebra
operation. Both pairs of graphs are supported on the same source
predicates. The resulting universal graph equivalences are
therefore precisely equality of the two pairs of morphisms in
$\mathcal B_T$. This proves uniqueness for every $P$. Taking $P$
to be terminal proves initiality in $\mathcal B_T^2$.
\end{proof}

\begin{proposition}[List initiality with parameters]
\label{prop:st-list-initiality}
For every object $V$ of $\mathcal B_T$, the endofunctor
$Z\mapsto1+V\times Z$ has an initial algebra stable under
parameters.
\end{proposition}
\begin{proof}
Use canonical codes of finite sequences and the predicate
\[
 \operatorname{List}_V(w)\quad\longleftrightarrow\quad
 w\text{ is a sequence code}\ \land\
       \forall i<\operatorname{len}(w)\,V(w_i).
\]
Lemma~\ref{lem:st-finite-sigma1} gives a $\Sigma_1$ presentation.
The empty sequence and the operation of prepending an entry are
primitive recursive, and $\mathrm{PA}$ proves their carrier
conditions and the unique nil-or-cons decomposition.

For an algebra in $\mathcal B_T/P$ on an arbitrary projection
$\xi:X\to P$, its nil operation is a morphism $P\to X$ over
$P$, and its cons operation has source
$(P\times V)\times_P X$. Evaluate a list from its final entry
back to its first entry, starting with the nil value in the fiber
over $p$. A finite sequence records these values and the
computation traces of the two algebra operations. As above, this
defines a $\Sigma_1$ relation. Induction in $T$ on length proves
its totality, functionality, and preservation of the projection to
$P$. Lemma~\ref{lem:syntactic-functional-search} makes it a morphism.
The nil and cons equations follow by splitting the record,
and induction in $T$ on length proves uniqueness of any morphism
satisfying these equations. This is initiality in every such slice.
\end{proof}

\subsection{Free variables and substitution}
\label{sec:st-substitution-realization}

The preceding propositions supply the initiality hypotheses of
Definition~\ref{def:categorical-syntax} for arbitrary fixed coefficient objects of bounded
arity in $\mathcal B_T$. The following algorithms realize its free-variable and substitution folds.

\begin{proposition}
\label{prop:st-substitution-realization}
On the preceding codes, the list operations, free-variable morphisms,
term substitution, and the option-valued formula-substitution morphism
of Appendix~\ref{app:syntax-theories} are represented by algorithms whose
defining equations and preservation of the specified carriers are
provable in $\PA$. Restricting each graph to its source predicate
gives a morphism of $\mathcal B_T$. These are the morphisms
constructed there by parameterized initiality. In particular, the
subobject $\Adm$ is the free-for substitution domain on
these raw named-variable codes, represented by a complemented
subobject in $\mathcal B_T$.
\end{proposition}
\begin{proof}
Concatenation, membership, and deletion on sequence codes are
primitive recursive. Membership compares numerical variable labels;
deletion removes every entry equal to the specified label.
Induction in $\mathrm{PA}$ on sequence length proves their nil and
cons equations and the two membership identities displayed in
Appendix~\ref{app:syntax-theories}. Each operation preserves its carrier whenever its
inputs satisfy their carrier predicates. Thus the corresponding
restricted graph equations define morphisms in $\mathcal B_T$.
The uniqueness assertion of
Proposition~\ref{prop:st-list-initiality} identifies them with the
appendix's list folds.

A postfix traversal computes the free-variable list. At a variable
it returns the singleton list; at function symbols, relation
symbols, and equality it concatenates the child lists in their
specified order; at negation it retains the child list; at
implication it concatenates the two lists; and at a quantifier it
deletes the bound variable. This is a primitive-recursive traversal
on finite codes. Induction in $\mathrm{PA}$ on node count proves
that its output satisfies $\operatorname{List}_V$ and proves the
constructor equations. These are the equations of the displayed
$K$-algebra on lists, so
Proposition~\ref{prop:st-syntax-initiality} identifies this algorithm
with the appendix's free-variable morphisms.

For substitution, use a traversal which retains at each node both
its original subtree and its computed result. Its parameters are
a variable code $x$ satisfying $V(x)$ and a term code $t$
satisfying $\operatorname{Tm}(t)$. Terms have a term result;
formulas have either a failure tag or a successful formula result.
At a term variable $v$, the result is $t$ if $v=x$ and is the
original variable term otherwise. At a function symbol, rebuild
the term from the substituted children. At an atomic formula,
rebuild the formula from the substituted terms and attach the
success tag. At negation and implication, use the lifted operations
on success-or-failure values.

At a quantified formula $\forall y\,\varphi$, let $z$ be the
computed option value for its body. The algorithm returns
\[
 \begin{cases}
  \operatorname{ok}(\forall y\,\varphi),&y=x,\\
  \widehat\forall(y,z),&y\ne x\text{ and }
       \bigl(x\notin\operatorname{FV}(\varphi)
             \text{ or }y\notin\operatorname{FV}(t)\bigr),\\
  \operatorname{fail},&\text{otherwise}.
 \end{cases}
\]
Here $x\notin\FV(\varphi)$ abbreviates failure of the numerical
list-membership test, and the other occurrence assertions have the
same meaning. The tests use numerical membership and equality. The first branch uses the original body, even if
the recursive computation for that body has failed. This is exactly
the quantified-constructor equation in Appendix~\ref{app:syntax-theories}.

The traversal is primitive recursive: it runs through the finite
postfix sequence, storing finite lists of original subtrees and
computed values, using only primitive-recursive coding operations.
Induction in $\mathrm{PA}$ proves that the term result is a valid
term and that every successful formula result is a valid formula.
The same induction proves that the first component is the original
subtree and verifies all the constructor equations of the
$K_{V\times\Tm}$-algebra used in Appendix~\ref{app:syntax-theories}.
By Proposition~\ref{prop:st-syntax-initiality}, its classes are
therefore precisely the appendix's $a_T$ and $b_F$.

Success and failure are decidable tags of this total option-valued
algorithm on its source. Their inverse images partition the source
in $\mathcal B_T$, giving the complemented subobject $\Adm$.
Write $\operatorname{Adm}(\varphi,x,t)$ for the numerical success
predicate. Induction in $\PA$ identifies it with the recursive
free-for condition. At a quantifier, this condition is
\[
 y=x\quad\text{or}\quad
 \bigl(y\ne x\ \land\
   (x\notin\operatorname{FV}(\varphi)
        \text{ or }y\notin\operatorname{FV}(t))\ \land\
   \operatorname{Adm}(\varphi,x,t)\bigr).
\]
The success value satisfies the usual substitution equation on
each such branch.

Finally, simultaneous induction in $\mathrm{PA}$ proves identity
substitution and substitution at an absent variable. At a
quantifier with $y=x$, the first branch returns the original
formula. With $y\ne x$, identity substitution passes the capture
test because $\operatorname{FV}(\eta x)=[x]$; absence of $x$
passes it because deletion of the distinct variable $y$ does not
remove $x$. The induction hypothesis supplies successful unchanged
results for the bodies in both cases. These are exactly the
source-restricted graph equations asserted at the end of
Appendix~\ref{app:substitution}.
\end{proof}


\subsection{Derivations and theories}
\label{subsec:syntactic-theories}

Let $\operatorname{fv}_{\Fm}(z)$ denote the numerical free-variable
list function just constructed. Since $0$ codes the empty list,
the predicate
\[
 \operatorname{Sent}(z)\equiv
 \operatorname{Fm}(z)\land\operatorname{fv}_{\Fm}(z)=0
\]
presents the sentence object
$S=(\omega,[\operatorname{Sent}]_T)$. Its inclusion into $\Fm$
is represented by the identity on sentence codes. This is the
pullback of the empty-list morphism along $\FV_{\Fm}$ required
by Definition~\ref{def:sentence-object}.

\begin{lemma}\label{lem:syntactic-derivation-initiality}
Let $T\supseteq\PA$, let $F$ be an object of $\mathcal B_T$, and
fix morphisms
\[
 p_n:R_n\longrightarrow F^n,\qquad
 q_n:R_n\longrightarrow F\qquad(0\leq n\leq N),
\]
where $N$ is a fixed nonnegative integer.
The inference endofunctor $H$ of
Appendix~\ref{app:closed-subobjects} has an initial algebra in
$\mathcal B_T/F$.  Its underlying object and conclusion morphism
are represented by $\Sigma_1$ formulas.
\end{lemma}
\begin{proof}
Choose $\Sigma_1$ formulas $F(z)$ and $R_n(r)$ for the displayed
objects and indices $u_n,v_n$ for $p_n,q_n$, respectively.
The nullary product is represented by the singleton code $0$.

First fix codes for finite ordered trees in which a node of arity
$n\leq N$ has a numerical label $r$.  Write
\[
 \operatorname{node}_n(r,d_1,\ldots,d_n)
\]
for the constructor.  These constructors can be chosen injective,
with pairwise disjoint images, and so that each $d_j$ is smaller
than the resulting code.  For example, use
$1+\langle n,\langle r,\operatorname{tuple}_n(d_1,\ldots,d_n)
\rangle\rangle$, with fixed-length tuple codes, taking the empty
tuple code to be $0$.  The pairing of Section~\ref{sec:syntactic}
is at least each of its components, so this inequality is provable
in $\PA$.
Course-of-values primitive recursion then defines a
primitive-recursive predicate $\operatorname{Tree}(d)$ and a
primitive-recursive list of the node occurrences of a tree.  At
each occurrence the list records the arity, numerical label, and
indices of its immediate children.  The list is ordered so that
children precede their parent.  Its constructor equations and
the unique decomposition of a tree are provable in $\PA$ by
induction on its numerical code.

Define $P(d)$ to assert that $d$ is such a tree and that, at every
node with label $r$, arity $n$, and immediate subtrees
$d_1,\ldots,d_n$, the following hold: $R_n(r)$, and there is an
$n$-tuple $(z_1,\ldots,z_n)$ such that
\[
 G_{u_n}(r,\langle z_1,\ldots,z_n\rangle)
 \quad\text{and}\quad
 G_{v_{m_j}}(r_j,z_j)\quad(1\leq j\leq n).
 \tag{\(*\)}
\]
Here $m_j$ and $r_j$ are the arity and label of the root of
$d_j$.  For $n=0$, the condition in \((*)\) is
$G_{u_0}(r,0)$.
Lemma~\ref{lem:st-finite-sigma1} makes $P$ a $\Sigma_1$
formula. The witnesses to node membership and computation belong
to its existential quantifiers; a derivation code records only
its rule labels and ordered children.

For a tree $d$ whose root has arity $n$ and label $r$, put
\[
 C(d,z)\ \equiv\ P(d)\land G_{v_n}(r,z).
\]
More precisely, this is the finite disjunction over $0\leq n\leq N$
using the primitive-recursive root projections.
The $T$-provable totality of $q_n$ on $R_n$ shows that
$T$ proves $P(d)\to\exists!z\,C(d,z)$ and
$C(d,z)\to F(z)$.
The formula $C$ is $\Sigma_1$ and its functionality is already
provable in $\PA$, by the functionality of the deterministic
program graphs.  Witness search therefore supplies an index
$e_c$ for which
\[
 \PA\vdash\forall d\forall z\,
             \bigl(G_{e_c}(d,z)\leftrightarrow C(d,z)\bigr).
\]
Its class represents a morphism $c:(\omega,[P]_T)\to F$ in
$\mathcal B_T$.

The constructor equations for the node list give, in $T$,
\begin{align}
 &P\bigl(\operatorname{node}_n(r,d_1,\ldots,d_n)\bigr)
                                                        \label{eq:syntactic-derivation-domain}\\
 &\quad\longleftrightarrow
 R_n(r)\land\bigwedge_{j=1}^{n}P(d_j)\land
 \exists z_1\cdots z_n\,
 \left(G_{u_n}(r,\langle z_1,\ldots,z_n\rangle)
       \land\bigwedge_{j=1}^{n}C(d_j,z_j)\right).
 \notag
\end{align}
For $n=0$, the last conjunct means $G_{u_0}(r,0)$.
The right-hand side is exactly a predicate presenting
$B_n(P,c)$.  Consequently the constructors define a total
morphism $\delta:H(P,c)\to(P,c)$ in $\mathcal B_T/F$:
its conclusion is $q_n(r)$ on this summand.

Let $\alpha:H(Q,q)\to(Q,q)$ be any algebra in this slice.
Represent $Q$ by $Q(y)$, $q$ by an index $b$, and $\alpha$ by
an index $a$.  We construct its fold without assuming that
$a$ is total on all numerical inputs.

Define $E_a(d,y)$ to assert that $d$ is a tree and there is a
finite list of values, one for each node occurrence, whose root
value is $y$, with the following property.  At a node of arity
$n$ and label $r$, if its children's values are
$y_1,\ldots,y_n$ and its own value is $y'$, then
\[
 G_a\bigl(\operatorname{tag}_n(r,y_1,\ldots,y_n),y'\bigr).
\]
The tag and tuple code here is the chosen presentation of the
$n$th summand of $H(Q,q)$.  This definition is $\Sigma_1$ by
the finite-sequence argument above.  The functionality of $G_a$
proves the functionality of $E_a$ in $\PA$: compare two value
lists in their children-before-parent order and use induction on
the position in the list.
Decomposing and concatenating value lists proves in $\PA$ that,
for trees $d_1,\ldots,d_n$,
\begin{align}
 &E_a\bigl(\operatorname{node}_n(r,d_1,\ldots,d_n),y\bigr)
                                                        \label{eq:syntactic-derivation-fold}\\
 &\quad\longleftrightarrow
 \exists y_1\cdots y_n\,
 \left(\bigwedge_{j=1}^{n}E_a(d_j,y_j)
       \land G_a(\operatorname{tag}_n(r,y_1,\ldots,y_n),y)
 \right).
 \notag
\end{align}
In particular, all witnesses on the right can be combined into
one evaluation list on the left; restricting a list gives the
reverse implication.

Induction on the numerical tree code now proves in $T$ that
each $d$ satisfying $P(d)$ has an evaluation $y$ in $Q$, and
that its image under $q$ equals its conclusion under $c$.
For the induction step, the induction hypotheses give values
$y_j$ for the immediate subtrees, with
$q(y_j)=c(d_j)$.  Equation~\eqref{eq:syntactic-derivation-domain}
therefore places
$\operatorname{tag}_n(r,y_1,\ldots,y_n)$ in the domain
$B_n(Q,q)$ of $\alpha$.  Its $T$-provable totality supplies
$y$, and its equation as a slice morphism supplies
$q(y)=q_n(r)=c(d)$.  Equation~\eqref{eq:syntactic-derivation-fold}
supplies the evaluation at $d$.  The same argument covers
$n=0$, with no induction hypotheses.
Thus $T$ proves
\[
 P(d)\longrightarrow\exists!y\,E_a(d,y),
 \qquad
 P(d)\land E_a(d,y)\longrightarrow Q(y),
\]
and, for all $d,y,z$, proves
\[
 P(d)\land E_a(d,y)
 \longrightarrow\bigl(G_b(y,z)\leftrightarrow C(d,z)\bigr).
\]
Witness search for the $\Sigma_1$ relation
$P(d)\land E_a(d,y)$ gives a partial-recursive graph uniformly
equivalent to it in $\PA$.  It defines a total slice morphism
$h:(P,c)\to(Q,q)$.  Equation~\eqref{eq:syntactic-derivation-fold}
is exactly the algebra-morphism equation
$h\delta=\alpha H(h)$.

If $g:(P,c)\to(Q,q)$ is another algebra morphism, its
algebra-morphism equation gives the same recursion equation
on the domain in~\eqref{eq:syntactic-derivation-domain}.
Induction on $d$, using the unique root decomposition, proves
in $T$ that its graph on $P(d)$ is equivalent to $E_a(d,y)$.
Outside $P$ both morphisms have empty domain, by their source
restriction equations.  Their graphs are therefore uniformly
$T$-provably equivalent.  Hence $g=h$, proving initiality.
\end{proof}

\begin{theorem}\label{thm:syntactic-generated-theory}
Let $T$ be a consistent extension of $\PA$ and put
$\mathcal B_T=\Tot(\Split(S_T))$. Fix a nonnegative integer $d$
and variable and signature objects $V,\Sigma_k,\Pi_k$ of
$\mathcal B_T$ for $0\leq k\leq d$.
The syntax, list, and substitution constructions above realize
the first-order structure of Appendix~\ref{app:syntax-theories}
in $\mathcal B_T$. For every axiom subobject $i_A:A\hookrightarrow S$
of its sentence object, the inference endofunctor $H_A$ has an
initial algebra with conclusion morphism $c_A:P_A\to\Fm$.
Its image is the least subobject of $\Fm$ containing the logical
and nonlogical axioms and closed under generalization and modus
ponens. The sentence pullback of this image is the least sentence
theory containing $A$.

Use the displayed sentence predicate, and choose the derivation
predicate $P_A(p)$ and conclusion index $e_A$ furnished by the
construction. Put
\[
 \operatorname{Th}_A(z)\equiv\exists p\,G_{e_A}(p,z).
\]
Then the restriction idempotent
\[
 t_A=[\operatorname{Sent}(z)\land\operatorname{Th}_A(z)]_T:
                  \omega\rightharpoonup\omega
\]
belongs to $S_T$, and its splitting represents that sentence theory.
\end{theorem}
\begin{proof}
Lemma~\ref{lem:syntactic-total-category} supplies regularity,
disjoint stable finite coproducts, and partial inverses for
monomorphisms. Numerical equality makes $V$ decidable.
Propositions~\ref{prop:st-syntax-initiality} and
\ref{prop:st-list-initiality} give the required initial algebras,
stable under every parameter object.
Proposition~\ref{prop:st-substitution-realization} identifies
their free-variable and substitution folds with the specified
numerical operations.

Every rule object in Appendix~\ref{app:first-order-derivations}
is a finite product, coproduct, or pullback of these objects and
$A$. Its premise and conclusion morphisms are composites of the
realized constructors and substitution morphisms. They are
therefore morphisms of $\mathcal B_T$.
Lemma~\ref{lem:syntactic-derivation-initiality} applies with $N=2$.
Theorem~\ref{thm:first-order-theory} and the sentence-theory
construction following it give the two leastness assertions.

The original morphism $[e_A]_T:\omega\rightharpoonup\omega$
has restriction $[P_A]_T$ and range $[\operatorname{Th}_A]_T$.
The graph
\[
 G_{e_A}(p,z)\land\operatorname{Sent}(z)
\]
is $\Sigma_1$ and PA-provably functional. By
Lemma~\ref{lem:syntactic-functional-search} it has a program
index $e_{A,S}$ with PA-provably equivalent graph.
It represents the sentence-conclusion pullback
$c_{A,S}:P_{A,S}\to S$. Its range in $S_T$ is precisely $t_A$.
By Lemma~\ref{lem:range-regular-image}, the same range in
$\Split(S_T)$ splits through the regular image of $c_{A,S}$.
Thus its splitting is the asserted sentence theory.
\end{proof}

For the arithmetic signature, a $\Sigma_1$ predicate
$\operatorname{Ax}_T(z)$ enumerating the axioms of a c.e.\ theory
$T$ defines the axiom subobject with predicate
$\operatorname{Sent}(z)\land\operatorname{Ax}_T(z)$.
Theorem~\ref{thm:syntactic-generated-theory} therefore represents
its internally generated sentence theory by a restriction
idempotent of $S_T$.
